% Pursuing Stacks (cotes 134-2 to 134-8): archive cote 134-5, site page 8 (PDF page 9), typescript page 195.
% typescript. Draft compared: ps4.tex:337-392.
% AI-transcribed by Claude Opus 5.5 (claude-opus-5-5, adaptive thinking, effort high), Message Batches API, batch ps2 (msgbatch_01NVwr1Yuo8qJf5rykeGBXPK), 2026-11-12; not yet checked by eye.
\scanpage{134-5/8}
\mspage{195}
description by generators and relations only when working simultaneously
with a whole bunch of base-points chosen with care -- or equivalently,
working in the algebraic context of \uline{groupoids} , rather than groups. Cho-%
sing paths for connecting the basepoints natural to the situation to \inserted{just} one
among them, and reducing the groupoid to a single group, will then hope-%
lessly destroy the structure and inner symmetries of the situation, and
result in a mess of generators and relations no one dares to write down,
because everyone feels they won't be of any use whatever, and just confu-%
se the picture rather than clarify it. I have known such perplexity my-%
self a long time ago, \inserted{namely in} \deleted{in \xout{\illegible{1}}} van Kampen type situations, whose
only understandable formulation is in terms of (amalgamated sums of)
groupoids. Still, standing habits of thought are very strong, and
during the long march through Galois theory, two years ago, it took me
weeks and months trying to formulate everything in terms of groups or
``exterior groups'' (i.e. groups ``up to inner automorphism'')\hand{,} and finally
\xout{\unsure{progressively}}
\xout{convincing myself only groupoids would fit nicely.} learning the lesson
and letting myself be convinced progressively, not to say reluctantly,
that \xout{only} groupoids only would fit nicely. Another ``technical point'' of
course is the basic fact(and the wealth of intuitions accompanying it)
that the Teichmüller groups are \xout{indeed} fundamental groups indeed -- a
fact ignored it seems by most geometers, because the \inserted[typed]{natural} ``spaces'' they are
fundamental groups of are not topological\xout{\illegible{1}} spaces, but \xout{multiplicities}
the modular ``multiplicities'' $M_{g,\hand{\nu}}$ - namely \uline{topoi}\hand{!} The ``points'' of these
``spaces'' are just the structures being investigated (namely algebraic
curves of type $(g,\hand{\nu})$), and the (finite) automorphism groups of these
``points''enter into the picture in a very crucial way. They can be \inserted{adequately} chosen
as part of the system of basic generators for the Teichmüller groupoid
$T_{g,\hand{\nu}}$ , \deleted{which} \inserted{The latter of course} is essentially (up to suitable restriction of base-points)
\inserted{just} the fundamental groupoid of $M_{g,\hand{\nu}}$. It is through this interpretation of
the Teichmüller groups or groupoïds that it becomes clear that the profi-%
nite Galois group $\mathrm{Gal}_{\overline{\underline{Q}}/\underline{Q}}$ operates on the profinite completion of these
and \inserted{of} their various variants, and \inserted{this} (it turns out) in a way respecting the
manifold structures and relationships tying them \inserted[typed]{tightly} together.

\margin{8}
